% Generated draft: requires mathematical and source review.
\documentclass[11pt]{article}
\usepackage[margin=25mm]{geometry}
\usepackage{fontspec}
\setmainfont{Latin Modern Roman}
\usepackage{amsmath,amssymb,mathtools,unicode-math}
\setmathfont{Latin Modern Math}
\usepackage{enumitem,needspace,xurl,hyperref}
\setlength{\parindent}{0pt}
\setlength{\parskip}{5pt}
\setlength{\emergencystretch}{3em}
\newcommand{\N}{\mathbb{N}}
\newcommand{\Z}{\mathbb{Z}}
\newcommand{\Q}{\mathbb{Q}}
\newcommand{\R}{\mathbb{R}}
\newcommand{\C}{\mathbb{C}}
\newcommand{\F}{\mathbb{F}}
\newcommand{\A}{\mathbb{A}}
\newcommand{\PP}{\mathbb{P}}
\newcommand{\E}{\mathbb{E}}
\newcommand{\eps}{\varepsilon}
\newcommand{\dd}{\,\mathrm{d}}
\begin{document}
\section*{UnsolvedMath ID: 30005737}
\textbf{Source problem code:} catalogue entry
\subsection{Definitions and mathematical statement}
For a finite-dimensional complex vector space \(V\), let \([\ ,\ ]\) and \(\{\ ,\ \}\) be Lie brackets on \(V\), defining \(\mathfrak g=(V,[\ ,\ ])\) and \(\mathfrak n=(V,\{\ ,\ \})\). A Lie algebra \(\mathfrak a\) is perfect if \([\mathfrak a,\mathfrak a]=\mathfrak a\), where the derived algebra is the linear span of its brackets. In the finite-dimensional setting, \(\mathfrak a\) is semisimple if its solvable radical, the largest solvable ideal, is zero. Here an ideal \(I\) is solvable if its derived series \(I^{(0)}=I\), \(I^{(k+1)}=[I^{(k)},I^{(k)}]\) reaches zero.

A post-Lie algebra structure on the ordered pair \((\mathfrak g,\mathfrak n)\) is a complex-bilinear product \(\cdot:V\times V\to V\) satisfying, for all \(x,y,z\in V\),
\[
\begin{aligned}
x\cdot y-y\cdot x&=[x,y]-\{x,y\},\\
[x,y]\cdot z&=x\cdot(y\cdot z)-y\cdot(x\cdot z),\\
x\cdot\{y,z\}&=\{x\cdot y,z\}+\{y,x\cdot z\}.
\end{aligned}
\]

\textbf{Conjecture, under the finite-dimensional interpretation.} For every finite-dimensional complex vector space \(V\) and every pair of Lie brackets on \(V\), if \(\mathfrak g\) is perfect and nonsemisimple and \(\mathfrak n\) is semisimple, then no such post-Lie product exists.

\textbf{Unresolved scope qualification.} The supplied statement does not explicitly impose finite dimensionality. The formulation above records the standard finite-dimensional interpretation as a provisional scope choice requiring human confirmation; it does not assert an extension to infinite-dimensional Lie algebras. The shared underlying vector space is made explicit to define the post-Lie identities. The supplied original excerpt refers to an unspecified preceding conclusion, so the nonexistence conclusion and the nonsemisimplicity hypothesis on \(\mathfrak g\) are taken from the supplied reconstructed statement.
\subsection{Short English statement}
The supplied conjecture says that a perfect, nonsemisimple complex Lie algebra cannot admit a post-Lie product with a semisimple second Lie bracket on the same vector space. The displayed formulation uses a provisional finite-dimensional interpretation because the supplied record leaves the dimensional scope unstated. The record labels the problem open and reports that the case \(\mathfrak n\cong\mathfrak{sl}_3(\C)\) is proved; these are supplied assessments requiring independent review.
\subsection{Sources}
\begin{itemize}
\item Local UnsolvedMath registry, permanent ID \(30005737\), source problem code \texttt{catalogue entry}, title \emph{Post-Lie Structures with Semisimple Target}; supplied collection: Oberwolfach Reports Open Problems. No individual proposer is identified; the supplied proposed year is 2023.
\item Supplied source citation: \emph{Mini-Workshop: Poisson and Poisson-type algebras} (2024), Oberwolfach Reports, pp.~2681--2715, DOI \texttt{10.4171/owr/2023/46}. Supplied literature-source label: \emph{Oberwolfach Report 46/2023}. \url{https://doi.org/10.4171/owr/2023/46}. The supplied original excerpt points to p.~2689 and conjectures an extension from \(\mathfrak{sl}_3(\C)\) to arbitrary semisimple \(\mathfrak n\).
\item Supplied literature assessment, checked on \texttt{2026-08-22}: the report is said to prove only the \(\mathfrak{sl}_3(\C)\) case and to present the general semisimple case as a conjecture; the supplied assessment says no later resolving proof was located. This is supplied dated curation, not an independently verified literature assessment. Its listed source is \url{https://doi.org/10.4171/owr/2023/46}.
\item The registry labels the statement \texttt{corrected\_verified} and says it was reconstructed and checked against the official report. These are supplied provenance claims; the report has not been independently consulted for this draft.
\end{itemize}
\end{document}
